\documentclass[script,con]{debate}
\motion{我厌蠢，我错了 / 我厌蠢，我没错}
\stance{反方}{我没错}
\meta{赛制}{中科大校赛}
\meta{口径来源}{备赛文档 第十节 反方口径表}
\begin{document}
\debatetitle
\begin{thesisbox}[全场一句话立论]常人的厌蠢，平常冲着教过了还犯：他被教过、提醒过，仍一再出错、不问不改，这样厌，看得有据，怪得公道；恼里那句「你本可以」，是对能负责的人提要求，这回放下，下回照样说、照样听。问完三句，这份厌平常的样子站得住。\end{thesisbox}

\begin{kit}{1. 反一开篇立论}{套件 · 组装后 802–825 字 / 180 秒}
\assembly{开头（任选 1） → 定义与判准（任选 1） → 立论主体}
\begin{pick}{开头}{1}
\begin{module}{如果质询拿到了对方承认凭记录的不争}
谢谢主席。刚才的质询里，对方已经认了：教过、被指出、又做错的，在看准了吗上不争。请评委先记下这一格，我方的立论就从这里开始。
\end{module}
\begin{module}{如果质询拿到了对方承认对那件事的恼火不反对}
谢谢主席。刚才的质询里，对方已经认了：对那件事本身的恼火，不反对。那今天要比的，就是这份恼平常停在事上，还是落到看低人。
\end{module}
\begin{module}{如果质询里对方回避了「他有没有被教过」}
谢谢主席。刚才的质询里，有一个问题对方没有正面回答：例子里的那个人，有没有被教过。请评委记住，这个问题我方全场还会再问。
\end{module}
\begin{fallback}{质询没有拿到明确成果}
谢谢主席。我先把我方的立场摆出来。今天这道题，请评委只盯一件事：常人说「我厌蠢」时，那份厌平常冲着什么，又往哪里走。
\end{fallback}
\end{pick}
\begin{pick}{定义与判准}{1}
\begin{module}{如果正方用了双方共用的中立定义与判准}
定义与判准，和对方大体一致；定义判准我们不争，我们争的是门槛。要改的如果是这份厌平常的样子，就是这份厌本身该改，哪怕改完还剩一点反感；要改的如果只在个别场合冒出来，那叫调分寸。平常就这样、只是改得掉，照样算错了。还有一处要说清：蠢，是本该会，却明显做错。第一次、没人教的不会，一般人不觉得他本该会；「明显」二字，排除鸡毛蒜皮。
\end{module}
\begin{module}{如果正方把看低写死在厌里}
对方把看低写死在厌里了。我方重述中立版：厌蠢，是觉得别人犯蠢时的嫌恶、反感，常带着想躲开的心；可以只冲着那件事，也可以冲着这个人、看低他；它是态度，不等于骂人或欺负人。蠢，是在一般人看来本该想到、做对的事上明显想错、做错，不论出于不会还是不肯。错没错，看清醒的我问完三句：看准了吗？对他公道吗？划得来吗？这份厌平常的样子该改，就是错了；只在个别场合调调分寸，就是没错。
\end{module}
\begin{fallback}{正方的定义与判准不在以上几种时}
定义与判准，我方用这一版：厌蠢，是觉得别人犯蠢时的嫌恶、反感，常带着想躲开的心；可以只冲着那件事，也可以冲着这个人、看低他；它是态度，不等于骂人或欺负人。蠢，是在一般人看来本该想到、做对的事上明显想错、做错，不论出于不会还是不肯。错没错，看清醒的我问完三句：看准了吗？对他公道吗？划得来吗？这份厌平常的样子该改，就是错了；只在个别场合调调分寸，就是没错。
\end{fallback}
\end{pick}
\begin{fixed}{立论主体}
我方要证明：问完三句，这份厌平常的样子站得住；对方要证明，它平常就该改。

我方第一点：教过了还犯。

常人的厌蠢，平常冲着教过了还犯：位置要他会的基本功，他被教过、提醒过，仍一再出错、不问不改。这样厌，\stress{看得有据，怪得公道}。

人的情绪，本来就分流。2004 年，一项汇总六十四项研究、一万二千多名被试的元分析发现：判断对方该负责，人就生气；判断他控制不了，人就同情。

本该，出自位置的公共要求，不是我会什么；本可以，看记录：他被教过没有、教过之后改不改。没有记录，先问；在问、在改还学不会的，给耐心。

每日商报 2024 年写过一位程序员：同事工作超过两年，还常把表格填错。他多次劝告，同样的错接二连三，关系一度降至冰点，现在很微妙。看准了吗？记录在手。对他公道吗？只怪他使得上劲的那一份。

我方第二点：恼他不弃他。

厌蠢收回的，是对他这一次的评价，作为人的尊重不动。Strawson 说，怨恨这类反应态度，是把对方当成要为自己负责的参与者；态度完全客观时，可以跟他谈判，却不能跟他讲理。说白了，恼里那句「你本可以」，是对能负责的人提要求；讲理由它带着，这是我们的推论，我方承认。

你想想《论语》里那一幕。子路当面说老师迂，孔子斥他一句「野哉」，接着把正名的道理讲完。朱熹的集注里，胡氏说，子路最终还是没领会。恼和讲理，在同一个当场；讲理，也不等他做对。

越线按时间判：这回放下，下回照样说、照样听，不越线；认定他说什么都没用，才越线，那是个别场合，要调分寸。这一点回扣对他公道吗。

我们全场会反复问一个问题：\stress{这个例子里，他有没有被教过？}没有，是我们先问、先教的那一种；有，正是我们守的那一种。

我们不说情绪无对错；我们说问完三句，那份平常冲着明显犯蠢的事去的反感，仍然站得住。谢谢。
\end{fixed}
\end{kit}
\begin{contingency}
\ifthen{正二问「看准了就够了吗」}{不够，所以我们同意三问都问；我们要证明的是，问完三句，这份厌平常的样子站得住。}
\ifthen{正二问「看准了这个人蠢，就能看低他吗」}{看低他的能力可以，拒绝给他作为人的尊重不行；我们的厌停在前一边。}
\ifthen{正方说我方辩护的是最好版本}{我们的限定不是前缀，是双方共用的定义；我们说的常人，会恼，会一时走开，下回会再说。这是常人，不是最好版本。}
\ifthen{时间不够}{先删每日商报一案的日期，再删孔子一段里朱熹那一句；定义与判准、两个论点的主张和后场问题不删。}
\end{contingency}

\begin{stage}{2. 反一被质询防守预案}{正二质询，120 秒双边计时}
\textbf{原则}：先答是或否，再用一句话守住前提，每答不超过十秒。被逼二选一时，指出问题的预设，给出我方表述。凡是推论，直接认是推论，再给依据。个别场合的毛病一律说「那一次是看错了」「那一步越线了，要调分寸」，不说「认错」。
\begin{qa}
\qarow{“也可以冲着这个人看低他”，是不是把鄙视塞进了厌？}{两头都写了，都用“可以”，没替任何一方选；平常落在哪一头、落在那里错不错，是今天要辩的。}
\qarow{按你们的定义，第一次不会坐高铁，算不算蠢？}{不算。定义里的一般人，是了解他处境、讲道理的一般人；第一次、没人教，这样的人不觉得他本该会。我们的定义不冤枉人。}
\qarow{那评论区为什么说“这不是谁都会吗”？}{评论区的多数票不是定义里的一般人，那是评论区看错了。他们那份也算厌蠢，看错是第一问要称的；但这是个别场合的误判，不是这份厌平常的样子。}
\qarow{看准了就够了吗？}{不够，所以我们同意三问都问；我们要证明的是，问完三句，这份厌平常的样子站得住。}
\qarow{看准了这个人蠢，就能看低他吗？}{看低他的能力可以，拒绝给他作为人的尊重不行；我们的厌停在前一边。}
\qarow{分寸调到不看低他，不就是改了你们平常的厌？}{看低不是这份厌平常的样子，只在个别场合冒出来，所以是调分寸。平常那份反感冲着事，它告诉我标准被踩了，我们守的就是它。}
\qarow{这份厌改得掉，按门槛规则就是错了吧？}{门槛不看改不改得掉，看平常还是个别。只在个别场合冒出来的，调分寸就好；我们证的是，平常那份冲着事的反感站得住。}
\qarow{厌了之后，是去交涉还是走开？}{厌，可以一时走开，但不把他划出讲理的范围：回头照样教、照样合作。}
\qarow{只在心里厌、从不表现，算不算厌蠢？}{算，心里那份就是厌蠢本身；它错不错，看判准。}
\qarow{新京报写同学同事，写的就是“我都可以做到你为什么做不到”，这不是比较吗？}{这几句我们全认。只凭“我都可以”、不看他学过没有的，是该先问的场合。新京报那段我们只取场景，记录由每日商报那一案补。}
\qarow{教爸妈付款，教三遍还按错，是不会还是不肯？}{看教过之后问不问、改不改：在问、在改还学不会的，给耐心；分不清时，先问他卡在哪。}
\qarow{会了的人估不准、自己看不见，你们凭什么说看得准？}{这些实验估的，都是一个你一无所知的人。被教过、劝过多次的同事不是不知情者，他学过没有，查得到，用不着估。没有记录的，我们也认该先问。}
\qarow{每日商报那案，劝了多次，接着就是“说了那么多遍，怎么还是那么蠢”？}{三句一起读：多次劝告，讲理没等他改好；冰点是一度，是降到过；现在是很微妙，不是断了往来。那句话没标说话人；之后讲不讲理，报道没写。}
\qarow{恼里混着轻蔑，你们承认吗？}{我们不说恼里没有轻蔑，说的是混着的那份往哪走。2007 年那项研究整句认：两种情绪常常一起出现，轻蔑可能从先前的生气里长出来；可能，不是平常的走向。往哪走，看回头还讲不讲。}
\qarow{恼里带着讲理，是不是你们的推论？}{是推论，立论里就说了。依据是 Strawson：态度完全客观，才不能讲理；客观态度装不下愤怒，混着恼的厌，到不了“完全”。}
\qarow{你们的“我”，是不是自我限定过的最好版本？}{我们的限定不是前缀，是双方共用的定义：一般人看来本该会、明显做错；第一次、没人教的，本来不算蠢。我们不靠谁的自我声明，靠记录。}
\end{qa}
\end{stage}

\begin{stage}{3. 反二质询正一问题链}{120 秒，双边计时}
\textbf{总目标}：让评委第一个回合就看到，正方的战场只剩两块：没有记录的比较，和看准之后还讲不讲理。拿三类承认：凭记录的、走神手滑的，在看准了吗上不争；对那件事的恼火不反对，一时走开不算越线；他们自己的例子里，那个人有没有被教过。对方说超过十五秒还没回到问题，礼貌打断；一条链拿不到就换下一条。拿到的承认，由反一在开篇第一句接住。
\begin{chain}{链一：让出的看准}{承认链：凭记录的、走神手滑的，在看准了吗上不争}
\q{教过、指出过、又做错，这种看准您方争不争？}{答“不争”：下一问；答“争”：那他被教过、劝过，还错在同一处，凭什么说看不准？}
\q{走神、手滑这类犯蠢，您方争不争看准？}{答“不争”：下一问；答“争”：走神手滑，也要先查他学没学过吗？}
\q{“一般人”是了解他处境的一般人吗？}{答“是”：收束；答“不是”：那第一次坐高铁不会，也算本该会？}
\cut{好，我理解为您方不争。下一个问题。}
\close{谢谢。凭记录的、走神手滑的，在看准了吗上，对方不争。}
\end{chain}
\begin{chain}{链二：恼火与越线}{承认链：对那件事的恼火不反对，一时走开不越线}
\q{对那件事本身的恼火，您方反不反对？}{答“不反对”：下一问；答“反对”：那是要人别有感觉、做圣人吗？}
\q{这回先不说、一时走开，算不算越线？}{答“不算”：下一问；答“算”：那越线看什么，看一时，还是看回头？}
\q{恼火不反对，去掉看低剩下那份，算不算错？}{答“不算”：收束；答“算”：剩下的就是恼火，您方刚才说不反对}
\cut{时间有限，我换个问法：一时走开，算不算越线？}
\close{谢谢。对那件事的恼火，对方不反对；一时走开，不算越线。}
\end{chain}
\begin{chain}{链三：例子分格}{边界链：把对方的例子分进「被教过」与「没被教过」两格；正一没举例子时，问他们准备用的例子}
\q{您方举的例子里，那个人有没有被教过？}{答“没有”：那是我方先问、先教的一种，下一问；答“有”：那正是我方守的一种，下一问}
\q{被教过、劝过多次的同事，您方的实验测过吗？}{答“没测过”：收束；答“测过”：请报研究名，我方驳论核对}
\cut{好，我理解为您方没有否认。}
\close{谢谢。他有没有被教过，是今天分歧的那一格；被教过的那一格，对方的实验没测过。}
\end{chain}
\end{stage}

\begin{kit}{4. 反二驳论}{套件 · 组装后 507–524 字 / 120 秒}
\assembly{开头 → 驳对方立论（任选 2） → 收尾}
\begin{fixed}{开头}
谢谢主席。我方想请评委把对方的立论放到判准下看。对方要证的，不是世上有一份厌得不对的厌，而是常人这份厌，平常的样子该改。我挑两处，一处一处看。
\end{fixed}
\begin{pick}{驳对方立论}{2}
\begin{module}{如果对方立论建立在「厌蠢平常凭比较，不凭记录」上}
对方说，厌蠢平常凭比较，不凭记录。可厌跟着代价走：2011 年的一组研究发现，事情与自己相关，推高的是愤怒，不推高别的情绪；轻蔑，在评判无能的行为时起作用。蠢的代价，反复落在一起做事的人身上：同一个组、同一间办公室，他被教过没有、改没改，本来就在手里。评论区里对陌生人的那句比较，我们认：算厌蠢，也看错了，是个别场合。2021 年一项对推特的研究还发现，网上的道德愤怒得到点赞，以后就更常这样表达。评论区，是被放大过的声音。
\end{module}
\begin{module}{如果对方立论建立在「会了的人估不准、自己看不见」上}
对方说，会了的人估不准，自己还看不见。这一类实验，我们不否认，只看估的是谁。后见之明的经典实验，被试估的是不知道结果的别人；替同龄人估字谜的那项，估的是一个不认识的同学，还被要求设想他和自己经历相同。估的都是一个他一无所知的人，偏差就长在这里。被教过、劝过多次的同事，不是不知情者：他学过没有，查得到，用不着估。要是开的药是先开口问，它只治得到没有记录的那一种；没有记录的，我们也认该先问。我们不评偏多少。
\end{module}
\begin{module}{如果对方立论建立在「劝了还犯，恼会长成轻蔑」上}
对方说，劝了还犯，恼会长成轻蔑。2007 年有一项研究，我们整句认：愤怒和轻蔑常常一起出现，轻蔑可能从先前的生气里长出来。可能，不是平常的走向；这些是预测因素，不是必要条件。同一篇的主干是另一句：生气当下冲上去，长远会和好；轻蔑才把人推出去。劝了不听，这回放下，下回照样说，就是还觉得话管得到他。理论分析里，轻蔑的特点之一，是不看好对方改善，恼里的「你本可以」，方向正好相反。真硬化成这个人说什么都没用，我们认越线，要调分寸。
\end{module}
\begin{module}{如果对方立论建立在「回头不讲理，越过了尊重」上}
对方说，回头就不讲理了，越过了尊重。这一步，有没有材料？每日商报那一案，三句要一起读：多次劝告，讲理没等他改好；关系一度降至冰点，是降到过；现在很微妙，不是断了往来。之后讲不讲理，报道没写。Strawson 划的线是：态度完全客观，才不能讲理；推不出没讲理就是完全客观，何况客观态度装不下愤怒。越线按时间划：这回先不劝、一时走开，不越线；他开口也不听、合作只下指令不给理由，越线，我们认要调。越线要证是平常，一则没写结尾的报道，证不了平常。
\end{module}
\begin{module}{如果对方立论建立在「厌蠢就是轻蔑、傲慢」上}
对方说，厌蠢就是轻蔑，是我在上、你在下。先看词：汉典里，厌是嫌恶、憎恶，看不起是鄙视的释义；中立定义两头都写了，平常落在哪头要证。再看理：平等的尊严，不排除按品质、成就区别评价人，只是要受承认尊重的约束。康德说，比较之下在心里看低别人，有时免不了，但表现出来就是冒犯；拒绝给人作为人应得的尊重，无论如何违背义务。同事同一处错了又错，我看低他在这件事上的用心，照样跟他讲理，这是区别评价，不是贬低人格。舆论批的傲慢，是表现出来的看低，我们不辩护。
\end{module}
\begin{fallback}{对方立论的主干不在以上几条时}
我方想把对方的立论放回判准。不管对方用什么论据，要证的都是同一件事：常人这份厌，平常的样子该改。证出有，不够；证出平常，才够。只证明这份厌有时看错、有时过头，也不够：那叫调分寸。存在更好的，也不等于这个就错了，替代态度只在证出平常的缺陷之后才登场。所以听对方每个例子时，问一句：这个例子里，他有没有被教过？没有，是我们先问、先教的那一种；有，正是我们守的那一种。教过、被指出、又做错这串记录，就是看准了吗最硬的依据。
\end{fallback}
\end{pick}
\begin{fixed}{收尾}
说到底，回到判准。分歧裁决写着：一方证出平常就有的东西，另一方只证出个别场合的东西，前者胜。有，不等于平常。我方的平常摆在桌上：教过了还犯，记录在手；恼他不弃他，回头还讲。你想想，对方哪一处证到了平常？
\end{fixed}
\end{kit}

\begin{stage}{5. 反四对辩要点}{90 秒，单边计时}
\textbf{目标}：B 块里把战场钉在一格：他有没有被教过。每次发言先一句回应，再抛一个问题；对方连续两次不答同一问，下一次发言点出来：「这个问题对方已经两次没有回答。」不开「划得来吗」的进攻。
\begin{points}{进攻点（4–5 个，每个 15–20 秒）}
\item 您方举的每个例子里，他有没有被教过？没有，是我们先问、先教的那一种；有，正是我们守的那一种。
\item 证平常的材料，除了评论区和一则报道，还有什么？评论区那句我们认：算厌蠢，看错了，个别场合。（对方的平常压在评论区时用）
\item 回头不讲理，报道写了吗？报道写的现状是很微妙，不是断了往来。这一步是材料，还是推论？
\item 对一个共事两年、劝过多次的同事，先开口问，还能问出什么？问了，他说会了又不改，归责落在谁身上？
\item 恼火您方不反对；去掉看低之后剩下的那份，算不算错？剩下的，正是我们守的那份。
\end{points}
\begin{points}{备用点（6–8 个）}
\item 对方说知道答案的人自己看不见：这些实验估的，都是一个他一无所知的人；同事学过没有，查得到，用不着估。
\item 对方说轻蔑可能从生气里长出来：可能，不是平常；预测因素，不是必要条件。同一篇说，生气长远会和好。
\item 对方说“说了那么多遍，怎么还是那么蠢”：报道没标说话人；“怎么还是”还在要求他，这是我们的推论。
\item 对方说对应偏差是人普遍的倾向：它咬的是处境没看见；教过、被指出、又做错，处境已经摆出来一回。
\item 对方说新京报写同学同事也是“我都可以”：这几句我们全认，只凭“我都可以”的，该先问；那段我们只取场景，记录由每日商报那一案补。
\item 对方说知识的诅咒：它是猜不准别人以为什么；对一个教过三遍的同事，用不着猜，记录在手。
\item 对方说尊重是底线：底线我们同意，我们本来没拿划得来吗去抵它；要证的是平常越线。
\item 对方说替代态度是先开口问：没有记录先问，本来就在我们的厌里；存在更好的，不等于这个就错了。
\end{points}
\begin{points}{对方最可能问的三个问题 → 回应}
\item 第一次做错那一下，手里有记录吗？→ 没有记录的，我们先问，不凭比较下判断；有人凭比较就厌了，那一次是看错了，要调分寸。
\item 看准了他这件事，就能看低他这个人吗？→ 看低他的能力可以，拒绝给他作为人的尊重不行；我们的厌停在前一边。
\item 先问、回头讲理，是常人平常就有，还是要学才有？→ 手里没记录，就还不知道他本该会没有，问是在拿记录，这是我们的推论；熟人之间，劝本身就是开口。
\end{points}
\end{stage}

\begin{stage}{6. 反三盘问问题链}{90 秒，单边计时}
\textbf{原则}：对方开始展开就打断：「谢谢，我听到了答案。」链一同一问先问正二、再问正四，两人不一致就是小结的第一句；链二先问正二，再连问正四两问。
\begin{chain}{链一：平常的材料}{问二辩 → 问四辩同一问题，检验口径}
\q{（问正二）除了评论区，常人平常凭比较，您方有材料吗？}{答“有”：请报机构、年份，转问正四；答“没有”：转问正四，看是否一致}
\q{（问正四）同一个问题：评论区以外，您方有材料吗？}{与正二一致：下一问；不一致：记下，小结重锤}
\q{（问正四）对共事两年、教过三遍的人，还要靠估吗？}{答“不用”：收束；答“要”：他学过没有，查不到吗？}
\close{谢谢两位。评论区以外的材料，我们记下了。}
\end{chain}
\begin{chain}{链二：看准之后}{钉推论与越线线}
\q{（问正二）回头不讲理，是报道写的，还是推论？}{答“推论”：转问正四；答“报道写的”：报道写的现状是很微妙，哪一句写了不讲理？}
\q{（问正四）这回先不劝、下回照样说，算不算越线？}{答“不算”：下一问；答“算”：那常人连劝都不能停一回？}
\q{（问正四）恼火不反对，去掉看低剩下的，算不算错？}{答“不算”：收束；答“算”：剩下的就是恼火，您方不反对的那份}
\close{谢谢两位。回头不讲理是推论；下回照样说，不越线。}
\end{chain}
\textbf{盘问目标}：小结里证明两件事。一，正方的平常，没有一起做事、记录在手的人的材料，或者两位辩手答得不一样，接回我方第一点与看准了吗；二，回头不讲理是推论，而这回放下、下回照样说不越线，接回我方第二点与尊重小问。
\end{stage}

\begin{stage}{7. 反二、反四被盘问防守预案}{两人口径必须一致}
\textbf{原则}：回答短；不确定时说「我方立论的立场是……」；两人的回答要能互相引用。被问是不是推论，直接认，再给依据。个别场合的毛病只说「那一次是看错了」「那一步越线了，要调分寸」。
\begin{qaa}
\qaarow{没有记录时，常人是先问，还是先判？}{先问。没有记录，就还不知道他本该会没有；有人凭比较就判了，那一次是看错了，要调分寸。}{先问。没记录先问，本来就在我们这份厌里；凭比较就判的，那一次是看错了，要调分寸。}
\qaarow{先问，是常人平常就有，还是要学才有？}{手里没记录，问是在拿记录，这是我们的推论；熟人之间，劝本身就是开口。}{问是在拿记录，这是我们的推论；熟人之间，劝本身就是开口。}
\qaarow{熟人第一次做错，你们手里有记录吗？}{第一次、没有记录的，我们先问，不凭比较下判断；记录攒在同一个组里，厌也跟着代价在那里起。}{没有，所以先问；我们恼的那一种，记录已经攒出来了。}
\qaarow{回头他开口先认定不值得听，算不算越线？}{算，越线，我们认要调分寸。那是个别场合；平常是这回放下，下回照样说、照样听。}{算，我们认要调。平常是这回放下，下回照样说、照样听。}
\qaarow{“怎么还是那么蠢”，说的是事还是人？}{报道没标说话人。“怎么还是”还在要求他，这是我们的推论；蠢字落在人身上，下调的是他这一方面的评价，平等尊严不排除区别评价。}{没标说话人；“怎么还是”还在要求他，是推论。下调的是他这一方面的评价，不是拒绝他作为人的尊重。}
\qaarow{恼里混着轻蔑，你们承认吗？}{我们不说恼里没有轻蔑，说的是混着的那份往哪走；回头还讲，就没走到那一头。}{不说没有，看往哪走。这回放下，下回照样说，就没走到那一头。}
\qaarow{研究说轻蔑可能从生气里长出来，你们怎么办？}{整句认：两种情绪常常一起出现，轻蔑可能从先前的生气里长出来。可能，不是平常的走向；预测因素，不是必要条件。}{整句认。可能不是平常；预测因素不是必要条件。}
\qaarow{下回还讲理，是常人自然会，还是要提醒自己？}{Strawson 说，正常人可以一时借客观态度躲开投入之累，但作为人，通常不能长久这样做。回来再说，就是常人的平常。}{一时躲开可以，不能长久，Strawson 说的是常人；回来再说，就是平常。}
\qaarow{新京报写同学同事，也是“我都可以做到你为什么做不到”？}{这几句我们全认：只凭“我都可以”、不看他学过没有的，该先问。那段我们只取场景，记录由每日商报那一案补。}{全认；只凭“我都可以”的该先问。新京报只取场景，记录看每日商报那一案。}
\qaarow{你们说的“我”，是不是自我限定过的最好版本？}{我们的限定不是前缀，是双方共用的定义；我们说的常人，会恼，会一时走开，下回会再说。这是常人，不是最好版本。}{不是。常人会恼、会一时走开、下回会再说；靠的是记录，不是自我声明。}
\qaarow{看准了这个人蠢，就能看低他吗？}{看低他的能力可以，拒绝给他作为人的尊重不行；我们的厌停在前一边。}{看低他这件事上的表现可以；拒绝他作为人应得的尊重，不行。}
\qaarow{康德说尊重的义务包括不持有轻蔑，心里那份也逃不过吧？}{心里那份也要称，我们不躲。百科说的轻蔑，是否认他有任何价值；我们看低的是他这件事上的表现，不否认他作为人的价值。}{心里那份也称。康德自己分了两件事：比较之下的看低有时免不了，拒绝基本尊重无论如何违背义务。}
\qaarow{对真正的笨，厌是错的吧？}{对努力了还学不会的人厌他，那一次是看错了，要调分寸；它不是常人厌蠢平常冲着的样子。}{在问、在改还学不会的，我们给耐心；有人厌了，那一次是看错了，要调分寸。}
\end{qaa}
\end{stage}

\begin{kit}{8. 反三盘问小结}{套件 · 组装后 508–538 字 / 120 秒}
\assembly{回应正三小结（任选 1） → 链一的结果（任选 1） → 链二的结果（任选 1） → 接回我方论点}
\begin{pick}{回应正三小结}{1}
\begin{module}{如果正三把战场引向「没有记录时凭比较」}
谢谢主席。正三刚才把战场引到没有记录的那一下。那一下我们认：没有记录，先问；有人凭比较就判了，那一次是看错了，要调分寸。这一格，我方从来不争。可门槛问的是平常，不是那一下。厌跟着代价走，代价反复落在一起做事的人身上，那里，他被教过没有，本来就在手里。
\end{module}
\begin{module}{如果正三把战场引向「怎么还是那么蠢」}
谢谢主席。正三刚才把战场引到那句「怎么还是那么蠢」。三句一起读：多次劝告，讲理没等他改好；关系一度降至冰点，是降到过；现在很微妙，不是断了往来。那句话报道没标说话人，之后讲不讲理，报道也没写。照字面读，「怎么还是」，还在要求他，这是我们的推论。
\end{module}
\begin{module}{如果正三把战场引向「先问是要学才有」}
谢谢主席。正三刚才说，先问要学才有，所以平常没有。可手里没有记录，就还不知道他本该会没有；问，是在拿记录，这是我们的推论。熟人之间，劝本身就是开口。有记录时，看记录；没记录时，先问。两样，本来都在我们这份厌里。所以先问，不是改出来的样子。
\end{module}
\begin{fallback}{正三小结的方向不在以上几条时}
谢谢主席。正三的小结，请评委用判准去量：对方要证的，是常人这份厌平常的样子该改。证出有，不够；证出平常，才够。一方证出平常就有的东西，另一方只证出个别场合的东西，前者胜。门槛不看改不改得掉，看平常还是个别。按这一条，我把盘问拿到的东西，一样一样摆出来。
\end{fallback}
\end{pick}
\begin{pick}{链一的结果}{1}
\begin{module}{如果对方承认平常的材料只有评论区和报道}
盘问第一条链，对方认了：平常凭比较，材料是评论区和报道。评论区那句，我们认：算厌蠢，也看错了，是个别场合。2021 年一项对推特的研究发现，网上的道德愤怒得到点赞，以后就更常这样表达。评论区，是被放大过的声音。对方的平常，压在一个没有代价、没有关系、也没有记录的地方。
\end{module}
\begin{module}{如果正二、正四答得不一样}
盘问第一条链，同一个问题，问了正二，又问了正四：评论区以外，常人平常凭比较，有没有材料。两位的答案不一样。各位评委可以记下来：守平常这一格的两个人，自己没有对齐。这一格，正是对方要守的平常。我方的平常一直是一句：厌跟着代价走，代价落在一起做事的人身上。
\end{module}
\begin{module}{如果对方回避了链一的问题}
盘问第一条链，评论区以外，常人平常凭比较有没有材料，这个问题对方没有正面回答。请评委记住。为什么难答？说有，就得报出来；说没有，平常就只剩评论区。而评论区那句我们认：算厌蠢，也看错了，是个别场合。厌跟着代价走，代价落在一起做事的人身上，那里记录现成。
\end{module}
\begin{fallback}{链一的结果不在以上几种时}
盘问第一条链，我替评委收一句：证平常，评论区不够。评论区那句我们认：算厌蠢，也看错了，是个别场合。也就是说，厌跟着代价走，代价反复落在一起做事的人身上：同一个组、同一间办公室，他被教过没有、改没改，本来就在手里。对方要证平常，得拿出那里的材料。
\end{fallback}
\end{pick}
\begin{pick}{链二的结果}{1}
\begin{module}{如果对方承认回头不讲理是推论}
盘问第二条链，对方认了：回头不讲理，是推论，报道没写。报道写的是：多次劝告，讲理没等他改好；关系一度降至冰点，是降到过；现在很微妙，不是断了往来。对方说的越线，压在一个推论上。按门槛规则，越线要证是平常；一个推论，证不出平常。这一条，各位评委可以记下。
\end{module}
\begin{module}{如果对方承认下回照样说不算越线}
盘问第二条链，对方认了：这回先不劝、下回照样说，不算越线。好，线就划在这里。常人的恼，劝了不听，这回放下，下回再说，就在线的这一边。越线，是他开口也不听、合作只下指令不给理由；那一步我们认要调分寸，可它要证是平常，不能只证可能。
\end{module}
\begin{fallback}{链二的结果不在以上几种时}
盘问第二条链，我替评委收一句：越线按时间划。这回先不劝、一时走开，不越线；他开口也不听、合作只下指令不给理由，越线，我们认要调分寸。每日商报那一案，多次劝告，关系一度降至冰点，现在很微妙；之后讲不讲理，报道没写。越线要证是平常。
\end{fallback}
\end{pick}
\begin{fixed}{接回我方论点}
说到底，盘问指向的，是我方两个论点，也是判准的两问。

第一点，教过了还犯，回扣看准了吗和归责。没有记录，先问；有记录，看记录。教过、被指出、又做错，还不问不改的，看得有据；只怪他使得上劲的那一份，怪得公道。

第二点，恼他不弃他，回扣对他公道吗里的尊重。恼里那句「你本可以」，是对能负责的人提要求。这回放下，下回照样说、照样听；认定他说什么都没用，那一步越线，要调分寸，那是个别场合。

说白了，有，不等于平常。你想想，自由辩里，我方还会问同一个问题：他有没有被教过？没有，我们先问、先教；有，正是我们守的那一种。
\end{fixed}
\end{kit}

\begin{stage}{9. 自由辩战场设计}{240 秒}
\begin{battleground}{战场一：平常在哪：评论区，还是一起做事的人（前 90 秒，主战场，回扣看准了吗）}
\begin{fields}
\field{目标}{让评委看到：厌跟着代价走，常人厌蠢平常面对的，是一起做事、记录在手的人；评论区那句我们认，是看错了的个别场合。}
\field{开场问题}{您方举的例子里，他有没有被教过？}
\field{对方答“没有”或“不知道” → 追问}{那是我们先问、先教的一种。要证平常，请拿出常人厌蠢时平常手里没有记录的材料，评论区以外的。}
\field{对方答“有” → 追问}{那正是我们守的一种：教过、被指出、又做错，这串记录，在看准了吗上就是最硬的依据。}
\field{对方可能反攻 → 我方防守}{“第一次做错那一下，只有比较。”→ 第一次、没有记录的，我们先问；有人凭比较厌了，那一次是看错了。“新京报写同学同事也是‘我都可以’。”→ 这几句全认，只凭“我都可以”的该先问；那段只取场景，记录由每日商报那一案补。“知道答案的人看不见偏差。”→ 那些实验估的是一个他一无所知的人；同事学过没有，查得到。}
\field{收束语}{厌跟着代价走；代价落在一起做事的人身上，那里记录现成。}
\end{fields}
\end{battleground}
\begin{battleground}{战场二：恼了之后，还讲不讲理（中间 90 秒，回扣对他公道吗里的尊重）}
\begin{fields}
\field{目标}{让评委看到：恼里那句“你本可以”是对能负责的人提要求；这回放下、下回照样说不越线；回头不讲理，对方只有推论。}
\field{开场问题}{回头不讲理，有直接材料吗，还是推论？}
\field{对方答“推论” → 追问}{报道写的现状是很微妙，不是断了往来；推论压不出平常。}
\field{对方答“有材料” → 追问}{请报机构、年份。每日商报那一案，三句一起读：多次劝告，一度降至冰点，现在很微妙；之后讲不讲理，没写。}
\field{对方可能反攻 → 我方防守}{“轻蔑可能从生气里长出来。”→ 整句认；可能不是平常，预测因素不是必要条件；同一篇说生气长远会和好。“怎么还是那么蠢，落在人身上。”→ 没标说话人；“怎么还是”还在要求他，这是我们的推论；下调的是他这一方面的评价。“孔子是圣人。”→ 我们只借这个样子：恼和讲理在同一个当场，讲理不等他做对；平常靠的是 Strawson 那条线。}
\field{收束语}{这回放下，下回照样说、照样听；恼他，不弃他。}
\end{fields}
\end{battleground}
\begin{battleground}{战场三：替代态度换得掉什么（最后 60 秒前插入，回扣门槛）}
\begin{fields}
\field{目标}{让评委看到：存在更好的，不等于这个就错了；对方的药，一半本来在我们的厌里，另一半管不到教过还不上心。}
\field{开场问题}{对劝过多次的同事，先开口问，还能问出新东西吗？}
\field{对方答“问他卡在哪” → 追问}{问了，他说不会，我们再教；他说会了、又不改，归责落在他那一份上，那正是我们恼的。}
\field{对方答“能” → 追问}{能问出来的，我们本来就问；没有记录先问，本来就在我们这份厌里。}
\field{对方可能反攻 → 我方防守}{“先问要学才有，就是平常没有。”→ 问是在拿记录，这是我们的推论；熟人之间，劝本身就是开口。“平常就这样、只是改得掉，照样算错。”→ 规则我们一字不改；要证的是看低是平常，不是改不改得掉。}
\field{收束语}{存在更好的，不等于这个就错了；对方的药，一半本来就在我们的厌里。}
\end{fields}
\end{battleground}
\begin{points}{必问（前两分钟内问完）}
\item 您方举的例子里，他有没有被教过？（全场问题）
\item 证平常的材料，除了评论区和一则报道，还有吗？
\item 回头不讲理，有直接材料吗，还是推论？
\item 恼火您方不反对，去掉看低之后剩下的那份，算不算错？
\end{points}
\begin{points}{必答（15 秒以内正面回答）}
\item 厌了之后，是去交涉还是走开？→ 可以一时走开，但不把他划出讲理的范围：回头照样教、照样合作。
\item 看准了这个人蠢，就能看低他吗？→ 看低他的能力可以，拒绝给他作为人的尊重不行；我们的厌停在前一边。
\item 只在心里厌，算不算？→ 算，心里那份就是厌蠢本身；它错不错，看判准。
\item 对真正的笨，厌是错的吧？→ 对努力了还学不会的人厌他，那一次是看错了，要调分寸；它不是常人厌蠢平常冲着的样子。
\end{points}
\textbf{时间分配与站位}：前 90 秒打战场一，一辩开场问全场问题，二辩接驳论打过的“估的都是陌生人”；中间 90 秒打战场二，三辩带盘问成果上（推论、越线线），四辩补尊重线；剩约 60 秒时插一轮战场三，由四辩收束，重复“存在更好的，不等于这个就错了；有，不等于平常”，不开新战场；留 20 到 30 秒备用。对方开新战场，与胜负无关的不接，一句“他被教过没有”带回战场一。一辩守定义与判准，二辩守论点一，三辩守论点二，四辩负责收束与价值。不主动开“划得来吗”的进攻。
\end{stage}

\begin{kit}{10. 反四结辩}{套件 · 组装后 890–918 字 / 210 秒}
\assembly{归纳分歧（任选 1） → 处理最强点（任选 1） → 封路（任选 1） → 论点收束 → 价值升华 → 结尾}
\begin{pick}{归纳分歧}{1}
\begin{module}{如果全场主战场落在「平常凭比较还是凭记录」}
谢谢主席。全场打得最多的，是平常凭比较还是凭记录。教过、被指出、又做错这串记录，我方一直按它判；没有记录的那一下，我们认先问，有人凭比较就判了，那一次是看错了。分歧只剩一格：常人厌蠢的时候，平常手里有没有记录。我方说有，因为厌跟着代价走，代价反复落在一起做事的人身上，那里记录现成。
\end{module}
\begin{module}{如果全场主战场落在「看准之后还讲不讲理」}
谢谢主席。全场打得最多的，是看准之后还讲不讲理。线按时间划：这回先不劝、一时走开，不越线；他开口也不听、合作只下指令不给理由，越线。分歧只剩一格：常人的恼，平常落在线的哪一边。我方说，落在这一边：这回放下，下回照样说、照样听。那句「你本可以」，本来就是还在要求他。越线我们认要调，那是个别场合。
\end{module}
\begin{fallback}{主战场不清楚时}
谢谢主席。全场打到最后，落在一处：常人说「我厌蠢」时，那份厌平常停在事上，还是落在看低人那一头。我方说，平常停在事上；看低和把人划出讲理的范围，只在个别场合冒出来。按分歧裁决，一方证出平常就有的东西，另一方只证出个别场合的东西，前者胜。评委要比的，是哪一边描述的平常，更像一个普通的、会说出这句话的人。
\end{fallback}
\end{pick}
\begin{pick}{处理对方最强点}{1}
\begin{module}{如果对方全场最有力的是「会了的人估不准、自己看不见」}
对方全场最有力的一点，是会了的人估不准、自己还看不见。我方不否认这些实验，只看估的是谁：估的都是一个他一无所知的人，不知道结果的别人，不认识的同学。偏差就长在一无所知的地方。被教过、劝过多次的同事，他学过没有，查得到，用不着估。方向固定，说的是偏的时候往哪边偏，回答不了偏的时候多不多；我们也不评偏多少。这把尺子，只在没有记录、只能靠它的时候，才是判断的全部依据。
\end{module}
\begin{module}{如果对方全场最有力的是「轻蔑可能从一次次恼火里长出来」}
对方全场最有力的一点，是轻蔑可能从一次次恼火里长出来。这篇研究我们整句认：两种情绪常常一起出现，轻蔑可能从先前的生气里长出来。三个词分清：可能，原文是 may，不是常常，更不是平常；预测因素，不是必要条件；一起出现，所以我们从不说恼里没有轻蔑。这篇的主干是另一句：生气长远会和好，轻蔑才把人推出去。往哪走，看回头：这回放下，下回照样说、照样听，就没走到那一头。真硬化成他说什么都没用，我们认越线。
\end{module}
\begin{module}{如果对方全场最有力的是「说了那么多遍，怎么还是那么蠢」}
对方全场最有力的一点，是那句「说了那么多遍，怎么还是那么蠢」。三句一起读：多次劝告，讲理没等他改好；关系一度降至冰点，是降到过；现在很微妙，不是断了往来。那句话，报道紧接着案例写下，没标说话人。照字面读，「说了那么多遍」是还在说，「怎么还是」还在要求他，这是我们的推论。蠢字落在人身上，下调的是对他这一方面的评价；平等的尊严，不排除区别评价。真到不再说、他开口也不听，我们认越线。
\end{module}
\begin{fallback}{对方最有力的一点不在以上几条时}
我方想请评委再听一遍我方最硬的一条。2004 年，一项汇总六十四项研究、一万二千多名被试的元分析发现：判断对方该负责，人就生气；判断他控制不了，人就同情，不同文化里结果高度一致。常人判断一份恼火对不对，本来就看可不可控、该不该他负责。有记录时，看准和归责一起过关：教过、被指出、又做错，还不问不改，恼得有据，也恼得公道。只凭一贯做不到，我们不厌；在问、在改还学不会的，给耐心。
\end{fallback}
\end{pick}
\begin{pick}{封路}{1}
\begin{module}{如果正方全场主打「平常就是凭比较」}
正方全场主打的，是平常就是凭比较。这条路的问题在材料：证平常，要证常人厌蠢时手里平常没有记录。待会儿的结辩，不管拿出哪一句话、哪一个例子，请评委用一个标准去检验：说话的人，跟被厌的人一起做过事吗？他被教过没有，说话的人知道吗？是对一无所知的陌生人，那一种我们认是看错了，是个别场合；是一起做事、记录在手的人，那正是我们守的。
\end{module}
\begin{module}{如果正方全场主打「就算看准了，也会滑到看低人」}
正方全场主打的，是就算看准了，恼也会滑到看低人。这条路的问题在最后一步：滑到那一步，是可能，还是平常？回头不讲理，有直接材料，还是推论？待会儿的结辩里，每一个滑过去的例子，请评委问两句：之后他开口，对面还听不听？这一步是写出来的，还是推出来的？推出来的，就请照推论来称。这回放下、下回照样说，在线的这一边。可能滑过去，不等于平常滑过去；越线要证是平常。
\end{module}
\begin{module}{如果正方全场主打「替代做得到，所以该改」}
正方全场主打的，是替代做得到，所以该改。这条路的问题在顺序：存在更好的，不等于这个就错了；替代态度，只在证出平常的缺陷之后才登场。待会儿的结辩，如果还是拿换一种态度来证错，请评委用两个标准检验：缺陷证在前了吗？那个替代，是不是本来就在我们的厌里？没有记录先问、对事批评、回头还讲，本来就是我们这份厌；纠正治的是不知道，治不了教过还不上心。
\end{module}
\begin{fallback}{正方全场路线不清楚时}
正方全场的路线，我不替他们归纳，只请评委带着判准去听待会儿的结辩。问三句：看准了吗，对他公道吗，划得来吗。再用门槛去称：要改的，是这份厌平常的样子，还是只在个别场合冒出来。一方证出平常就有的东西，另一方只证出个别场合的东西，前者胜。证出有，不够；证出平常，才够。存在更好的，也不等于这个就错了。最后说的每一句，都请放到这杆秤上。
\end{fallback}
\end{pick}
\begin{fixed}{论点收束}
我方的两个论点，打到现在，还站着。

第一点，教过了还犯，回扣看准了吗和归责。情绪本来就按该不该他负责分流；该不该他负责，看记录，记录查得到。没有记录先问，新来的先教，在问在改还学不会的给耐心；教过还不问不改的，看得有据，怪得公道。

第二点，恼他不弃他，回扣对他公道吗里的尊重。厌蠢收回的，是对他这一次的评价，作为人的尊重不动。态度完全客观，才不能讲理；说白了，恼里那句「你本可以」，是对能负责的人提要求。孔子斥完子路，照样把道理讲完。一时走开不越线；认定他说什么都没用，才越线，那是个别场合，我们认要调。
\end{fixed}
\begin{fixed}{价值升华}
最后，你想想这样一个人。

不是你的老师，是带过你的那位前辈。同一张表，你第三次交错了。他没替你改，也没说算了。他把表推回来：「这个上回讲过。你本可以的。再来一遍，我看着。」

那一下，你脸上发烫，心里有点恼他。

可你回头想。最该怕的，不是他恼。是哪天他不恼了：接过表，自己改完，笑一笑，什么也不说。那一刻，你会觉得，他已经不指望你了。

我们今天争的，不是该不该烦。是烦的时候，\stress{还把不把对面那个人，当成一个能做对的人}。

我没错，不是我比你强。是我把你当成能负责的人：你本可以，所以我恼。在意标准，也在意一起做事的人；为在意这件事，我不必认错。
\end{fixed}
\begin{fixed}{结尾}
所以，恼他，是因为还指望他成，还把他当成能负责的人。我厌蠢，我没错。谢谢。
\end{fixed}
\end{kit}
\begin{contingency}
\ifthen{时间不够}{先删处理最强点模块里的研究细节，再把封路压成检验标准那一句；论点收束、价值升华与结尾不删。}
\ifthen{正方在自由辩之后发动奇袭质询，追问“平常停在事上有什么材料”}{反四在归纳分歧之后正面答一句：平常靠机制，厌跟着代价走，代价落在一起做事的人身上，那里记录现成；评论区那句我们认是看错了。不展开，直接进处理最强点。}
\ifthen{结辩前对方提出全场没出现过的新论据}{封路组用兜底模块，请评委用判准去听，不追着新论据驳。}
\end{contingency}

\begin{stage}{11. 反方奇袭预案}{质询 150 秒双边 / 申论 120 秒，二选一}
\textbf{执行人}：奇袭质询由反三执行，奇袭申论由反一执行；执行人在环节结束时举手向主席示意发动。不在开场就用掉，越晚越清楚对方哪里最疼；进入反四结辩就没有机会了。首选窗口四。
\begin{fields}
\field{窗口一：正一立论之后}{反二质询紧接着就是取证，一般不发动。只有正一把看低写死在厌里、正二质询时又坚持这一版，才在正二质询反一之后发动奇袭质询，用链四钉定义。执行人：反三。正方用中立定义时不发动。}
\field{窗口二：对辩之后}{发动条件：正四在对辩里说出“回头不讲理有材料”“先问要学才有”这类可以钉死的表态。发动奇袭质询，用链二或链三取证。执行人：反三。}
\field{窗口三：盘问小结之后、自由辩之前}{发动条件：我方论点二在盘问小结里被打成“恼会长成轻蔑”，或“恼是看好他改，还是不看好他”还没有人正面打。发动奇袭申论，立这个战场，给自由辩定议题。执行人：反一。}
\field{窗口四：自由辩之后、结辩之前}{发动条件：对方全场把平常压在评论区和一则报道上，或回头不讲理始终没有直接材料。发动奇袭质询，用链一或链二，逼正四在结辩里回答。执行人：反三。首选窗口。对方没有可以钉死的表态时，改用奇袭申论，执行人：反一。}
\end{fields}
\begin{chain}{奇袭质询链一：如果对方全场把平常压在评论区和一则报道上}{逼出「平常没有一起做事的人的材料」}
\q{平常除了评论区和一则报道，还有材料吗？}{答“没有”：下一问；答“有”：请报机构、年份、写的是谁}
\q{那一则报道，写了回头不讲理吗？}{答“没写”：下一问；答“写了”：报道写的现状是很微妙，哪一句写了？}
\q{评论区里的人，跟被厌的人一起做过事吗？}{答“没有”：收束；答“有”：那他被教过没有，说话的人知道吗？}
\cut{好，我理解为您方没有别的材料。下一个问题。}
\close{谢谢。对方的平常，压在评论区和一则没写结尾的报道上。}
\end{chain}
\begin{chain}{奇袭质询链二：如果对方的回头不讲理始终没有直接材料}{逼出「是推论」}
\q{回头不讲理，有直接材料吗，还是推论？}{答“推论”：下一问；答“有材料”：请报出处，我方结辩核对}
\q{报道写的现状，是很微妙，还是断了往来？}{答“很微妙”：下一问；答“断了往来”：请念原文}
\q{一个推论，能证出平常吗？}{答“不能”：收束；答“能”：那门槛规则里的平常，还要材料吗？}
\cut{时间有限，我换个问法：这一步是写出来的，还是推出来的？}
\close{谢谢。回头不讲理是推论；报道写的是很微妙。}
\end{chain}
\begin{chain}{奇袭质询链三：如果对方把替代态度定为先开口问}{逼出「替代一半在我们的厌里」}
\q{劝过多次的同事，先开口问，还能问出新东西吗？}{答“不能”：下一问；答“能”：能问出来的，我们本来就问}
\q{问了，他说会了又不改，归责在不在他？}{答“在”：下一问；答“不在”：那他使得上劲的那一份算谁的？}
\q{没有记录先问，算不算在我们的厌里？}{答“算”：收束；答“不算”：我方一辩第一点就说了没有记录先问}
\close{谢谢。对方的药，一半本来就在我们的厌里，另一半管不到教过还不上心。}
\end{chain}
\begin{chain}{奇袭质询链四：如果对方把看低写死在厌里}{逼出「看低是可以，不是必然」}
\q{汉典里，看不起是厌的释义，还是鄙视的释义？}{答“鄙视”：下一问；答“厌”：请念汉典原文}
\q{中立定义两头都写了“可以”，您方认吗？}{答“认”：下一问；答“不认”：那您方用的是哪一版定义？}
\q{只冲着那件事的反感，算不算厌蠢？}{答“算”：收束；答“不算”：那评论区里那句“我厌蠢症犯了”冲着的是什么？}
\close{谢谢。看低是可以，不是必然；平常落在哪头，要证。}
\end{chain}

\begin{kit}{奇袭申论}{套件 · 组装后 513–516 字 / 120 秒}
\assembly{开头 → 为什么现在发动（任选 1） → 立一个战场 → 收尾}
\begin{fixed}{开头}
谢谢主席。我方用掉这一次奇袭，只做一件事：立一个立论里没有展开的战场，请评委听完再判。恼里那句「你本可以」，传过去的到底是什么。
\end{fixed}
\begin{pick}{为什么现在发动}{1}
\begin{module}{如果在盘问小结之后发动}
盘问小结刚刚结束，自由辩还没开始。我方想先替自由辩定一个议题，请双方在接下来的自由辩里正面打：恼，是看好他能改，还是不看好他。我方先把答案摆出来。这个问题，决定恼他和弃他，是不是一回事。
\end{module}
\begin{module}{如果在自由辩之后发动}
自由辩刚刚结束，结辩还没开始。我方想在最后两篇发言之前，把一个问题交给对方四辩，请他在结辩里正面回答：恼，是看好他能改，还是不看好他。我方先把答案摆出来。这个问题，决定恼他和弃他，是不是一回事。
\end{module}
\begin{fallback}{发动时机不在以上几个窗口时}
我方选在这里发动，是因为有一个问题，双方都要在剩下的环节里回答：恼，是看好他能改，还是不看好他。双方都绕不开它，请评委替双方盯着。我方先把答案摆出来。这个问题，决定恼他和弃他，是不是一回事。
\end{fallback}
\end{pick}
\begin{fixed}{立一个战场}
先看恼传过去的是什么。1984 年一项实验里，孩子做砸之后，大人冲他生气，孩子更多想的是「我没用心」；大人表示同情，孩子反倒更多觉得「我能力不行」。那是儿童实验，搬到成年人身上，是我们的推论。可方向值得听：没用心，是下回改得了的事。

再看轻蔑。理论分析里，轻蔑的特点之一，是不看好对方改善。恼里那句「你本可以」，方向正好相反：还在说你本可以，就是还看好你改，这一步也是推论。

还有一项 2006 年的经理研究，测的是经理信不信人能改。它不测厌，可它说明了一件事：决定教不教的，是人能不能改这个看法。恼里的「你本可以」，站的正是相信他改得了的那一边，这是我们的推论。

最后，Roberts 概括 Bell 的区分：怨恨针对行为、促使交涉；轻蔑针对人、促使退出。我们的恼冲着这件事、回头还讲，在怨恨这一格。

你想想，恼他和弃他，方向是反的。恼，是还指望他成。
\end{fixed}
\begin{fixed}{收尾}
说到底，还是判准的第二问：对他公道吗。恼里那句「你本可以」，是把他当成能负责的人，作为人的尊重不动。看好他改，就没有弃他。请对方回答：恼，是看好他能改，还是不看好他？
\end{fixed}
\end{kit}
\textbf{不发动的情况}：场上没有出现可以钉死的表态，只是我方预设的矛盾；我方两个论点都没有被打塌，自由辩的战场已经在盘问小结里立起来；对方已经在自由辩之后发动奇袭质询，反四要把时间留给结辩；反四结辩已经开始。宁可不用，也不要把奇袭申论念成第二篇立论。
\end{stage}
\end{document}
